\section{Exercise 3: Folding}
\subsection{Task and method}
Implement \texttt{[yn, yorigin] = fold(xn, xorigin)} for
\begin{equation}
  y(n)=x(-n).
\end{equation}
Folding reflects the signal about $n=0$. Each sample at time $n_i$
moves to $-n_i$ while retaining its amplitude. To keep the output vector
ordered by increasing time, its elements must be stored in reverse order.

Let $N$ be the input vector length. Input position $i$ becomes output
position $j=N+1-i$. The corresponding time relation is
\begin{align}
  j-\texttt{yorigin} &= -(i-\texttt{xorigin}),\\
  N+1-i-\texttt{yorigin} &= -i+\texttt{xorigin},\\
  \texttt{yorigin} &= N+1-\texttt{xorigin}.
\end{align}
In Scilab, \texttt{xn(\$:-1:1)} selects the elements from the last
position to the first, using a step of $-1$. Both the vector order and
the origin parameter must be updated.

\subsection{Source code}
\lstinputlisting[caption={Scilab folding function and example call},label={lst:fold}]{../3.sci}

\subsection{Example and expected result}
Use the same input sequence as in the preceding exercises:
\begin{lstlisting}[numbers=none]
[yn, yorigin] = fold([1, -2, 3, 6], 3);
\end{lstlisting}
The calculated output is
\[
  \texttt{yn}=[6,3,-2,1],\qquad
  \texttt{yorigin}=4+1-3=2.
\]
\begin{table}[H]
  \centering
  \caption{Time reversal of each input sample}
  \label{tab:fold-mapping}
  \begin{tabular}{lrrrr}
    \toprule
    Sample value & 1 & $-2$ & 3 & 6\\
    \midrule
    Input time $n_i$ & $-2$ & $-1$ & 0 & 1\\
    Reflected time $-n_i$ & 2 & 1 & 0 & $-1$\\
    \bottomrule
  \end{tabular}
\end{table}
Reading the reflected samples in increasing time gives
\begin{table}[H]
  \centering
  \caption{Output vector arranged in increasing time}
  \label{tab:fold-output}
  \begin{tabular}{lrrrr}
    \toprule
    Output vector position & 1 & 2 & 3 & 4\\
    \midrule
    Output time $n$ & $-1$ & 0 & 1 & 2\\
    Output value $y(n)$ & 6 & 3 & $-2$ & 1\\
    \bottomrule
  \end{tabular}
\end{table}

\subsection{Execution evidence}
\resultfigure{images/exercise3.png}
  {Input and folded output for Exercise 3}
  {fig:fold}
The captured Scilab plot agrees with the calculated output coordinates
in \Cref{tab:fold-output}.

\subsection{Discussion}
The captured output is the reflection of the input about the vertical
axis at $n=0$. The sample of amplitude $6$ moves from $n=1$ to $n=-1$,
and the sample of amplitude $1$ moves from $n=-2$ to $n=2$. The sample
at $n=0$ retains its amplitude $3$, although its position within the
stored vector changes from three to two. Folding changes time
coordinates, not the signs of the amplitudes.
